% Generated draft: requires mathematical and source review.
\documentclass[11pt]{article}
\usepackage[margin=25mm]{geometry}
\usepackage{fontspec}
\setmainfont{Latin Modern Roman}
\usepackage{amsmath,amssymb,mathtools,unicode-math}
\setmathfont{Latin Modern Math}
\usepackage{enumitem,needspace,xurl,hyperref}
\setlength{\parindent}{0pt}
\setlength{\parskip}{5pt}
\setlength{\emergencystretch}{3em}
\newcommand{\N}{\mathbb{N}}
\newcommand{\Z}{\mathbb{Z}}
\newcommand{\Q}{\mathbb{Q}}
\newcommand{\R}{\mathbb{R}}
\newcommand{\C}{\mathbb{C}}
\newcommand{\F}{\mathbb{F}}
\newcommand{\A}{\mathbb{A}}
\newcommand{\PP}{\mathbb{P}}
\newcommand{\E}{\mathbb{E}}
\newcommand{\eps}{\varepsilon}
\newcommand{\dd}{\,\mathrm{d}}
\begin{document}
\section*{UnsolvedMath ID: 8800015}
\textbf{Source problem code:} catalogue entry
\subsection{Definitions and mathematical statement}
Let $X$ be a variety over a finite field $\F_q$, $q=p^a$. A standard computational goal is its zeta function
\[
Z(X,T)=\exp\left(\sum_{r\geq1}\#X(\F_{q^r})\frac{T^r}{r}\right).
\]
In a suitable $p$-adic cohomology theory with compact supports this is encoded by Frobenius through
\[
Z(X,T)=\prod_i\det(1-TF_q\mid H_c^i(X))^{(-1)^{i+1}},
\]
with the cohomology and Frobenius normalization fixed so the trace formula applies.

The source asks to develop a higher-dimensional analogue of a particular nondegenerate-curve Frobenius-lift method. Such an analogue would need to specify the variety class, lift, cohomological model, finite basis/reduction procedures, precision control, and effective Frobenius computation. A general zeta-function algorithm is not automatically the requested analogue of that specific method.

One common toric nondegeneracy convention for a Laurent polynomial $f$ with Newton polytope $\Delta$ requires that, for every face $\tau$, the face polynomial $f_\tau$ and all logarithmic derivatives $x_j\partial f_\tau/\partial x_j$ have no common zero in the algebraic torus. This is explanatory context, not a supplied restriction to hypersurfaces: the record does not identify its nondegeneracy definition, dimensions, toric compactifications or complete-intersection extensions, nor reproduce the curve method.

No input encoding, characteristic restrictions, required output precision, or complexity bound is given. Accordingly the target is a method-development programme; the supplied reports of higher-dimensional cohomology and surface algorithms do not establish a fully general analogue with unprovided guarantees.
\subsection{Short English statement}
Extend the source's nondegenerate-curve Frobenius-lift algorithm to higher-dimensional varieties, with effective cohomological reduction and precision control. The exact curve method and intended variety/complexity class are missing.
\subsection{Sources}
\begin{itemize}
\item Local UnsolvedMath registry, permanent ID 8800015, exact source code catalogue entry, ``Higher-dimensional nondegenerate Frobenius algorithms''; attributed to Kiran Kedlaya, no proposed year supplied. Registry status partially solved, classification PARTIAL-PROGRESS.
\item \emph{IPAM Workshop on Number Theory and Cryptography---Open Problems} (2006), attachment 6220, PDF pages 30--32, logical slide 16, \url{https://www.ipam.ucla.edu/programs/workshops/workshop-i-number-theory-and-cryptography-open-problems/?tab=schedule}. Supplied pdf-and-legacy-ppt visual-audit extraction accessed July 29, 2026. The source audit covered 17 PDFs and two legacy PPT decks and excluded repeated animations/solved definitions; no PDF/PPT was accessed here.
\item The supplied summary names Lauder--Wan, Harvey and Kedlaya surveys for higher-dimensional zeta algorithms, and ``Castryck--Chtcherbakov, Gajovic, and others'' for toric/surface work, without precise titles or theorem hypotheses. It reports polynomial dependence on $p$, degree and $\log_pq$ for some algorithms, but supplies no dimension or input-model conditions; no blanket polynomial-time statement is inferred.
\item The summary says important higher-dimensional cases, especially surfaces, are implemented while general nondegenerate methods remain open. Provenance flags current primary-literature status and redistribution terms NEEDS\_REVIEW.
\end{itemize}
\end{document}
